\reportheading{3.\quad 任务 1：顶棚射流验证与门洞影响}
\subhead{3.1\quad 与实测温度剖面的对照}
47.2~kW 时，OP-1 粗网格计算的平均温升比实测低 192~K。OP-6 的平均误差仅为 $+6$~K（RMSE 为 33~K），但实测峰值在燃烧器上方，计算峰位向左偏了 0.40~m。106.2~kW 时，OP-6 的剖面误差最大。实测与计算峰值均位于封闭侧，分别在燃烧器左侧 0.20 与 0.80~m。
\newcommand{\validationpanel}[3]{%
\begin{tikzpicture}
\begin{axis}[width=.48\linewidth,height=3.70cm,xmin=-340,xmax=340,ymin=0,ymax=1350,
xtick={-300,0,300},ytick={0,600,1200},scaled ticks=false,
axis lines*=left,axis line style={ink!55},tick style={ink!55},
tick align=outside,major tick length=2pt,ymajorgrids,grid style={ink!13},
title={#2},title style={at={(axis description cs:0,1)},anchor=south west,font=\sffamily\small\bfseries,color=ink,inner xsep=0pt},
tick label style={font=\sffamily\scriptsize},enlarge x limits=false]
\addplot[only marks,mark=*,mark size=1.15pt,ink] table[col sep=comma,x=figure_r_cm,y=exp_rise_K]{../../05_Calculations/#1_90_120.csv};
\addplot[teal,line width=1.05pt,no marks] table[col sep=comma,x expr=\thisrow{model_x_cm}-330,y=fds_rise_K]{../../05_Calculations/#1_90_120.csv};
\node[anchor=north east,font=\sffamily\scriptsize,text=black!70,inner sep=2pt] at (rel axis cs:1,1) {RMSE #3 K};
\end{axis}\end{tikzpicture}}
\reportfigure{{\sffamily\small 温升 (K)\hfill
\textcolor{teal}{\rule[.4ex]{14pt}{1.1pt}} FDS（粗网格、绝热）\quad
\textcolor{black!75}{$\bullet$} Peng 等\par}
\begin{tabular}{@{}c@{\hspace{4pt}}c@{}}
\validationpanel{op1_47_coarse}{OP-1\enspace 47.2 kW}{192}&\validationpanel{op1_106_coarse}{OP-1\enspace 106.2 kW}{82}\\
\validationpanel{op6_47_coarse}{OP-6\enspace 47.2 kW}{33}&\validationpanel{op6_106_coarse}{OP-6\enspace 106.2 kW}{335}
\end{tabular}\par{\sffamily\small 距火源轴线距离 (cm)}}{图 1.\quad 粗网格上 90--120~s 的平均顶棚温升与 文献~\cite{peng} 图 9 的数字化实验数据。图 2 另对高功率 OP-6 的壁面边界进行对照。}

Peng 等~\cite{peng}给出的测温相对不确定度最高为 8\%。以 OP-1、70.8~kW 为例，其图 2 的纵轴为温度，图 9(c) 为温升，两图的数值尺度也有差异。为统一比较依据，四个工况均取图 9 所绘温升。实验数据取自稳定阶段，模拟采用 90--120~s 的均值。将这段记录前后两个 15~s 的均值逐点比较，四个工况各自最大的测点变化在 9.2 至 28.2~K 之间。
\subhead{3.2\quad 网格与模型参数敏感性}
在 OP-1 低功率工况中，网格逐级加密后，RMSE 从 192.4 降至 187.7，再降至 176.2~K。高功率下对应数值为 82.0、75.7 和 79.9~K，中网格误差最小。对高功率粗网格，换用 WALE 后，RMSE 从 Deardorff 的 82.0 降至 77.0~K。保留 Deardorff 并将辐射份额从 0.29 调到 0.20，RMSE 为 74.6~K，调到 0.50 则达 131.1~K。就这个工况而言，辐射份额对温升预测的影响更大。

\clearpage
\subhead{3.3\quad 门态对换气的影响与壁面散热敏感性}
在 106.2~kW 下，仅保留 R1 开启后，门洞总进气量从 OP-1 的 0.318~kg/s 降至 OP-6 的 0.095~kg/s。R1 门口由进气转为排气的位置也随之降低，从 0.30--0.38~m 移至 0.22--0.30~m（图 B1）。在 47.2~kW 下比较同一高度，OP-6 从 0.38~m 处排出的气流达 212.1$\,^{\circ}$C，OP-1 为 64.8$\,^{\circ}$C。R1 低位进入的空气可能将羽流推向封闭侧，方向与实测峰位偏移一致。

对 106.2~kW 的 OP-6，保持网格、火源和门态不变，用 15~mm 导热板代替绝热壁后，峰位误差由 0.70 降至 0.10~m，整条剖面的 RMSE 从 361.4 降至 100.7~K。在 R1 一侧，距燃烧器 1.0--3.3~m 的 24 个测点仍平均低估 112.5~K（图 2、表 E1）。

导热板工况的气相能量收支覆盖房间及门外延伸域。90--120~s 内，净辐射（\texttt{Q\_RADI}）为 $-61.4$~kW，气流带走的显焓（\texttt{Q\_CONV}）为 $-34.2$~kW。FDS 给出的 \texttt{Q\_COND} 为 $-9.6$~kW，包含气体与表面的对流换热以及开放边界的热扩散~\cite{fdsuser}。板内导热则通过改变板面温度，影响板面与气体的换热。
\reportfigure{\begin{tikzpicture}
\begin{axis}[name=profile,width=.90\linewidth,height=3.8cm,xmin=-340,xmax=340,
ymin=0,ymax=1380,ytick={0,600,1200},xtick={-300,0,300},xticklabels=\empty,
ylabel={温升 (K)},ymajorgrids,legend columns=3,legend style={at={(.5,1.02)},anchor=south}]
\addplot[ink,only marks,mark=*,mark size=1.0pt] table[col sep=comma,x=figure_r_cm,y=exp_rise_K]{../../05_Calculations/op6_106_board_doorflow_90_120.csv};\addlegendentry{Peng 等}
\addplot[rust,dashed,thick] table[col sep=comma,x expr=\thisrow{model_x_cm}-330,y=fds_rise_K]{../../05_Calculations/op6_106_doorflow_90_120.csv};\addlegendentry{绝热壁}
\addplot[teal,thick] table[col sep=comma,x expr=\thisrow{model_x_cm}-330,y=fds_rise_K]{../../05_Calculations/op6_106_board_doorflow_90_120.csv};\addlegendentry{替代板材}
\end{axis}
\begin{axis}[at={(profile.south west)},anchor=north west,yshift=-.24cm,
width=.90\linewidth,height=3.0cm,xmin=-340,xmax=340,ymin=-200,ymax=200,ytick={-150,0,150},xtick={-300,0,300},xlabel={距火源轴线距离 (cm)},ylabel={误差 (K)},ymajorgrids]
\path[fill=teal!10] (axis cs:100,-200) rectangle (axis cs:330,200);
\addplot[ink!65,dashed] coordinates {(-330,0) (330,0)};
\addplot[teal,thick] table[col sep=comma,x expr=\thisrow{model_x_cm}-330,y=error_K]{../../05_Calculations/op6_106_board_doorflow_90_120.csv};
\node[font=\sffamily\scriptsize,text=ink] at (axis cs:210,125) {平均 $-112.5$ K};
\end{axis}\end{tikzpicture}}{图 2.\quad 106.2~kW、OP-6 的顶棚温升。两条 FDS 曲线采用相同诊断网格和 90--120~s 时间窗。下图为板材模型逐点预测值减去文献~\cite{peng} 的实验值的残差。}

在 47.2~kW 的 OP-1 粗网格上，同一块导热板使 RMSE 从 192.4 增至 243.3~K。最热点仍在燃烧器上方，但温升从 435 降至 340~K，实测为 580~K。进一步检查绝热工况，同位置的气体平均仅比热电偶测珠高 4.7~K（图 E1）。另在门流诊断网格上保持分辨率不变，将外边界从 $y=1.40$ 延至 1.85~m，RMSE 从 166.6 变为 173.8~K。两次检查中，顶棚温升仍然偏低。任务 2 在其自身的绝热壁面和排烟条件下单独评估净高。
